% ==============================================================================
% banq-front.tex — Unified front matter for the bundle.
% Single title page, single arXiv tag, combined executive abstract, master TOC.
% Loaded by 000.banq-all/banq.tex in bundle mode only.
% ==============================================================================

% Page 1: title block, arXiv tag, and huge centred bank emoji.
\thispagestyle{empty}
\maketitle

% Single arXiv-style vertical tag (one date for the whole bundle).
% ==============================================================================
% arxiv-tag.tex
% The arXiv-style vertical preprint tag printed on the first page of each
% standalone paper and on the bundle title page. The date is taken from
% \banqArxivDate (defined in banq-meta.tex) and may be overridden locally
% before % ==============================================================================
% arxiv-tag.tex
% The arXiv-style vertical preprint tag printed on the first page of each
% standalone paper and on the bundle title page. The date is taken from
% \banqArxivDate (defined in banq-meta.tex) and may be overridden locally
% before % ==============================================================================
% arxiv-tag.tex
% The arXiv-style vertical preprint tag printed on the first page of each
% standalone paper and on the bundle title page. The date is taken from
% \banqArxivDate (defined in banq-meta.tex) and may be overridden locally
% before \input{../shared/arxiv-tag} via \renewcommand.
% ==============================================================================

\begin{tikzpicture}[remember picture, overlay]
  \node[rotate=90, anchor=center, font=\Huge, color=codegray]
    at ([xshift=1cm]current page.west)
    {\banqArxivId~~\banqArxivDate};
\end{tikzpicture}
 via \renewcommand.
% ==============================================================================

\begin{tikzpicture}[remember picture, overlay]
  \node[rotate=90, anchor=center, font=\Huge, color=codegray]
    at ([xshift=1cm]current page.west)
    {\banqArxivId~~\banqArxivDate};
\end{tikzpicture}
 via \renewcommand.
% ==============================================================================

\begin{tikzpicture}[remember picture, overlay]
  \node[rotate=90, anchor=center, font=\Huge, color=codegray]
    at ([xshift=1cm]current page.west)
    {\banqArxivId~~\banqArxivDate};
\end{tikzpicture}


\vspace*{\fill}
\begin{center}
  \twemoji[height=8cm]{bank}
\end{center}
\vspace*{\fill}
\clearpage

% Page 2: combined executive abstract and keywords. Styled to match the
% per-paper verso abstracts emitted by \banqPaper (large bold "Abstract"
% heading, \small body, vertically centred via \vspace*{\fill} pads).
\thispagestyle{empty}
\vspace*{\fill}
\begin{center}{\large\bfseries Abstract}\end{center}
\medskip
\begingroup
\small
\noindent
DeFi lending markets channel billions of dollars but share a structural weakness: a sharp price move triggers liquidations, the dumped collateral pushes prices lower, and a cascade follows. XPower Banq is a permissionless lending protocol designed to dampen these cascades — along with several other systemic risks — without giving up the openness and composability that make DeFi work.

\medskip
\noindent
This volume gathers the protocol's complete documentation in one bound document. It comprises six companion papers reproduced here without abridgement — five of them individually published as arXiv-style preprints — organised into five parts.

\medskip
\noindent
\textbf{Part~I — Protocol.} The whitepaper introduces five design choices that set the protocol apart from incumbents such as Compound and Aave~\cite{compound2019,aave2020}. (1)~\emph{Optionally locked positions:} borrowers may lock their position for a fixed term in exchange for liquidation immunity; the locked share $\phi$ attenuates cascade severity by a factor of $(1{-}\phi)$. (2)~\emph{Lethargic governance:} every parameter change is bounded multiplicatively ($0.5\times$--$2\times$ per cycle) and phased in asymptotically, so users always have time to react. (3)~\emph{Beta-distributed position caps} slow how fast a single account can scale up, yielding Sybil resistance that grows as $\sqrt{n{+}2}$. (4)~\emph{Transferable debt positions} with inverted ERC20 semantics let liquidators \emph{assume} bad debt rather than repay it — a far more capital-efficient model. (5)~\emph{Log-space TWAP oracles} compute time-weighted geometric means in both directions, cheaply and without overflow.

\medskip
\noindent
\textbf{Part~II — Engineering Primitives.} Two implementation papers describe the lower-level on-chain building blocks the protocol depends on. A 16-slot quarterly \emph{ring-buffer time lock}~\cite{lck} tracks total locked depth in $O(1)$ via the algebraic identity $D = \Sigma Q - Tt + pL$, avoiding the gas cost of iterating slots. A \emph{log-space compounding interest index}~\cite{log} stores the cumulative sum $L = \sum r_i$ in place of the multiplicative product, eliminating overflow and replacing write-path exponentiation with addition.

\medskip
\noindent
\textbf{Part~III — Theory.} Formal foundations and proofs covering continuous compounding, log-space oracle aggregation, time-weighted parameter integration, and token-bucket rate limiting; proofs that the protocol resists liquidation cascades, Sybil amplification, abrupt governance, and MEV front-running; and a Nash-equilibrium analysis of lock adoption that identifies two self-reinforcing equilibria — low and high — with utilization-dependent dynamics.

\medskip
\noindent
\textbf{Part~IV — Evaluation.} Four empirical studies validate the protocol. \emph{Capacity accumulation} under the beta cap confirms the predicted $O(\sqrt{n})$ Sybil-resistance scaling. \emph{Liquidation-cascade simulations} across five market-impact regimes show up to 80\% cascade reduction at full lock adoption. \emph{Log-space TWAP oracles} are stress-tested under price shocks and liquidity changes across 60 on-chain Foundry scenarios, exhibiting two-tick manipulation immunity. \emph{Bad-debt risk} is quantified by Merton jump-diffusion Monte Carlo and corroborated by a closed-form analytical bound: at the default 66.67\% LTV bad debt remains bounded under empirical crash distributions, and the 33\% conservative-mode floor keeps bad debt at zero for crashes up to 50\%.

\medskip
\noindent
\textbf{Part~V — Reference.} A consolidated glossary of over~120 terms — protocol mechanisms, mathematical notation, and DeFi terminology — intended to be read alongside the body of the volume.

\medskip
\noindent
The defaults strike a deliberate balance between capital efficiency and solvency: 66.67\% LTV with a 50\% over-collateralisation buffer, adjustable through lethargic governance (\S\ref{subsec:lethargic-governance}).

\medskip
\noindent\textbf{Keywords:} DeFi, lending protocol, liquidation, oracle, TWAP, governance, interest rates, time locks, log-space arithmetic, Sybil resistance, Nash equilibrium\par
\endgroup
\vspace*{\fill}
\clearpage

% Master table of contents — single hierarchical listing covering all parts and
% the chapters within them. Children's \section calls feed into this TOC.
% Force the TOC to start on the next odd (recto) page; see \banqForceRecto.
\banqForceRecto
{
  \hypersetup{linkcolor=primaryblue}
  % Compact 2-column master TOC. The "Contents" heading is emitted at full
  % size above the columns; the listing itself is rendered inside multicols
  % at \small with tightened vertical spacing and narrower number boxes so
  % the bundle's deep numbering ("IV.\,4.6.1") still fits in narrow columns.
  \section*{\contentsname}%
  \begingroup
    \small
    \setlength{\columnsep}{1.5em}%
    \setlength{\cftbeforesecskip}{2pt}%
    \setlength{\cftbeforesubsecskip}{0.5pt}%
    \setlength{\cftbeforesubsubsecskip}{0pt}%
    \setlength{\cftbeforeparaskip}{0pt}%
    \setlength{\cftsecnumwidth}{3.0em}%
    \setlength{\cftsubsecnumwidth}{3.7em}%
    \setlength{\cftsubsubsecnumwidth}{4.6em}%
    \begin{multicols}{2}
      \makeatletter\@starttoc{toc}\makeatother
    \end{multicols}
  \endgroup
}
\bigskip
\hrule
\bigskip
% Force the TOC itself to occupy an even number of pages, so Part~I starts on
% an odd (recto) page in duplex bindings; see \banqForceRecto.
\banqForceRecto
